% !PRINTED 315
\weberpage{315}
%<<CONT
Theiler haben wir schon betrachtet (\S.\,70 bis 72), und alle echten Theiler von $S$ sind metacyklisch. Darunter sind Theiler vom $21^{\mathrm{sten}}$ und $7^{\mathrm{ten}}$ Grade, die wir mit $T_{21}$, $T_7$ bezeichnen wollen, ferner Theiler vom $24^{\mathrm{sten}}$ Grade, und noch andere Theiler, deren Gradzahlen in 24 aufgehen, die wir hier alle in dem Zeichen $T_{24}$ zusammenfassen wollen.

Diese Gruppen $T_{24}$ können nicht die sieben Grössen $(\xi) = (\xi_1, \xi_2, \xi_3)$ transitiv mit einander verbinden, weil der Grad einer transitiven Permutationsgruppe von sieben Elementen durch 7 theilbar sein muss (Bd.\,I, \S.\,154, 7.). Wir wollen nachweisen, dass die den Gruppen $T_{24}\,Q$ entsprechenden Permutationsgruppen der $X$ imprimitiv sind. % [sic] Full stop before the syntactic continuation with \glqq und dass\grqq{}.
und dass diese Gruppen daher von unserer Betrachtung ausscheiden\footnote{Aus der Literatur über die Gleichungen $8^{\mathrm{ten}}$ Grades ist zu erwähnen: \gesperrt{Nöther}, Mathem. Annalen, Bd. XV. \gesperrt{Wilshaus}, Ueber die algebraische Auflösbarkeit der Gleichungen $8^{\mathrm{ten}}$ Grades. Dissertation. Marburg 1888.}.

Da alle Permutationen $\sigma$, die zwei der sieben Elemente $(\xi)$ ungeändert lassen oder nur unter einander permutiren, ein drittes Element, nämlich die Ergänzung zum Tripel, ungeändert lassen, so sind zwei Fälle möglich:

\begin{itemize}
\item[1)] die Gruppe $T_{24}$ lässt ein Element in Ruhe, oder
\item[2)] die sieben Elemente werden in zwei Theile von drei und vier Elementen zerlegt, von denen jeder Theil nur unter sich permutirt wird.
\end{itemize}

Im letzten Falle können wir die drei Elemente als einem Tripel, und demnach die vier Elemente als einem Quadrupel angehörig betrachten. Denn werden drei Elemente, die nicht einem Tripel angehören, unter sich permutirt, so werden auch die drei Ergänzungen je zweier Elemente zum Tripel unter sich permutirt, und das eine übrig bleibende Element bleibt ungeändert. Wir kommen also auf den Fall 1) zurück.

Wir fügen nun zu den sieben Systemen $(\xi)$ noch das achte $(0,0,0) = (0)$ hinzu und bezeichnen die acht Grössen $X_{\xi_1, \xi_2, \xi_3}$ mit

\weq{1}{0, 1, 2, 3, 4, 5, 6, 7.}

Im Falle 1) bleiben dann durch die Substitutionen von $T_{24}$ zwei dieser acht Elemente, etwa 0, 1, ungeändert. Es giebt ferner in $Q$ eine und nur eine Substitution $\gamma_0$, durch die 0 mit 1 vertauscht wird.
